\section{Conclusion}\label{sec:conclusion}

We have described \ITP, a proof assistant built to be the environment of a learning system: a dependent type theory with inductive families, an impredicative bottom universe, quotients and a structural guard, decided by a C++ kernel of about \locKernelHeader{} lines; a Miller-pattern unifier and certificate-checked decision procedures for polynomial arithmetic; a second, independent checker; and a layer that turns the assistant into an environment, with a finite menu of actions and a typed term builder for the rest, probes that try an action and take it back, work counted in kernel steps, exact hindsight, a bank of \bankProblems{} problems, and a Gymnasium interface with no text in the loop.

What the evidence supports is modest and concrete.
A probe of an action costs a median of \probeMed\,ms and a decision of the Gymnasium interface \decMed\,ms on one core of a two-core machine; the kernel checks the median reference proof of the bank in \chkKernelMedianMs\,ms, so that the verdict can be had for every finished proof, and the independent checker confirms it about \chkRatioMedian{} times more slowly, as a witness.
The soundness of the kernel is tested, not proved, by a frozen reference kernel run in lockstep, a second checker, canaries, a fuzzer, \mutTotal{} mutants, sanitizers and a log of \findNumbers{} findings.
The testing also reached its own limits: writing the rules down for this report turned up two places where the kernel checks less than the usual conditions on a motive and on the positivity of a family (\cref{app:deviations}); the second checker, written from the same specification, shares the first omission in part and does not examine the second.
The second omission can be exploited: with excluded middle and propositional extensionality added to a world, a family that the kernel admits gives a proof of $\sy{False}$ that both checkers accept (\cref{app:d2}).

What the evidence does not show is whether a learner profits.
The report has no learning experiment, and its claims about the design are claims about an interface: a verdict the proposer cannot influence, states that can be restored, actions that are well formed by construction, a cost that is a function of the task, and failures that yield theorems.
Whether these properties help is an empirical question that needs a fixed protocol, an untrained control, and a budget measured in the unit that the environment provides; the environment is meant to make that experiment cheap to run and its outcome checkable.

\paragraph{Future work.}
The two deviations of \cref{app:deviations} are not repaired in this version and should be, in the kernel, with a gate for each (for the second, the proof of $\sy{False}$ of \cref{lst:dev-exploit} must fail at the declaration of \code{R}), and the independent checker should re-check the admission of inductive families.
The remaining gaps of the calculus, universe polymorphism, mutual and nested families and a real-number library, bound the mathematics that the bank can contain; the bank itself needs a second reader of its statements and proofs for its hard tier.
A faster implementation of the policy interface would move the throughput of a decision toward that of the kernel.
Part of the metatheory, starting with the guard and the forcing of indices, could be mechanised.
And the environment should be used: to train and to compare learners on the bank and on synthetic theorems, against an untrained control, with the budgets of this report.

\subsection*{Availability and reproducibility}\label{sec:availability}

The code, the dataset of the problem bank, the sources of this report and the scripts that produced its numbers are in one repository at commit \mCommit{} (\reporturl); a snapshot is archived with the report (\reportdoi), under the licence \artifactlicense.
The directory \code{preprint/} of the repository holds the sources of the report (\code{src/}), the bibliography, the measurement scripts, the listings and their outputs, and the raw measurements as JSON and logs; \code{python3 -I preprint/build.py} regenerates every macro that holds a number, and then the PDF, from them.
The scripts and their costs on the machine of \cref{sec:measure} are: \code{measure.py}, the library chain, its re-check by the kernel and by the independent checker, and the micro-benchmarks (about three minutes); \code{measure\_check.py}, the cost of trusting a sample of reference proofs (about ten minutes, 2\,GB); \code{src/python/experiments/exp\_bench.py}, the latency and the throughput of the environment (a few minutes after the world is built); \code{measure\_assembler.py}, the round trip of the policy interface; \code{measure\_gates.py}, the gates named in \cref{tab:testing} and the sanitizer run; and \code{run\_listings.py}, the listings of \cref{sec:usage} and \cref{app:deviations} (those in Lean need a Lean~4).
The code, the dataset and the world are identified by digests so that a rerun can be compared with ours: the digest of the directory of the kernel's sources (\code{src/c++/hybrid}) is \code{\kernelDigest}, that of the names and types of the \libGlobals{} globals of the library world is \code{\worldDigestShort\ldots}, and that of the outcomes of the benchmark tasks is \code{\detDigest}.
The counts of this report (lines, declarations, mutants, findings, problems, proofs) are read from the files of the repository and do not depend on the machine; the times do.
